
\section{Ocean products (ORB)}
\label{sec:orb_product}

\subsection{Product Description}
\label{sec:orb_prod_description}
The \acs{orb} product, either as defined in \cite{CEOS-ARD_PFS_ORB} or in the variants outlined in the proposal for this study, can be relatively straightforwardly derived from the \acs{nrb}. This derivation provides both the necessary full resolution and spatially averaged \gls{sigma0} values on a map-projected grid. 

\acs{orb} comprises multiple layers of \gls{sigma0}, each corresponds to a specific layer of preprocessing. These preprocessing steps include denoise, outlier removal, downsampling and surface roughness normalisation. The raw (full resolution) and intermediate products are also saved, including the noise-equivalent \gls{sigma0} (NESZ), a bright target detection mask (full resolution), and the normalization factor used for surface roughness correction. Comprehensive annotations detailing the processing steps applied to each layer ensures full traceability and allows users to reverse or reproduce the processing chain as needed.

A full resolution product is useful for the detection of ships or other anthropogenic targets, which could be done starting from the NRB product. Providing a \gls{sigma0} product is a logical choice when treating the data as an image of a locally homogeneous distributed surface scatterer. It also is worth noting that a full-resolution \acs{orb} product is, in general, not a good starting point for the generation of enhanced ocean products. In many cases, the generation of these products requires SLCs (or even raw data), for example for sub-look generation to enable the computation of intensity cross-spectra, or for the optimal suppression of azimuth ambiguities in coastal areas by optimal selection of the azimuth-wavenumber (also known as Doppler frequency) weighting. 

For applications where the NRCS is relevant (e.g. surface stress or wind, oil-spill detection), a low-speckle spatially averaged product including polarimetric information, and a compensation of the known part of the angle of incidence dependence is preferable to a full-resolution product which first needs to be properly despeckled, noting that speckle filtering requires a level of expertise that makes the product deviate from the ARD rationale. Moreover, given the limitation of the effective resolution by the azimuth cut-off, which is largely wavelength independent, a spatially averaged product to a generally achievable spatial resolution facilitates the creation of intercomparable products by different sensors.

% The \acs{ceos} \acs{ard} Ocean Radar Backscatter (ORB) product is a full-resolution intensity product with the necessary corrections applied to allow a nominal interpretation as a Normalized Radar Cross Section (NRCS, $\sigma_0$). The \acs{ceos} specification of the ORB product, closely related to the NRB product, includes geoid correction, to convert to $\sigma_0$ but it does not include any despeckling. This choice is questionable given that the azimuth resolution of a SAR image is limited not by the radar characteristics or the processed Doppler bandwidth, as over land, but by the coherence time of the surface as observed by the radar. This results in the so-called azimuth cut-off,
% \begin{equation}
%         \lambda_c \approx \frac{R\cdot \sigma_v}{v_\mathrm{orb}},
%     \end{equation}
% with $R$ the slant-range to the surface and $v_\mathrm{orb}$ the orbital velocity of the satellite. The result is that the effective azimuth resolution of the intensity image is limited to \qtyrange{50}{400}{\meter}, depending on the sea state, which, in turn, is mostly wind-driven. Most notably, this effective resolution is largely wavelength independent, which would, in principle, facilitate the creation of intercomparable products by different sensors, if these are first averaged to a generally achievable spatial resolution.

% A full resolution product is still useful for the detection of ships or other anthropogenic targets, for which the $\sigma_0$ convention is not particularly useful. Ship detection could be done starting from the NRB product. Providing a $\sigma_0$ product is a logical choice when treating the data as an image of a locally homogeneous distributed surface scatterer. %Given the limitation of the effective resolution by the azimuth cut-off, it is surprising that the ORB product is defined as a full nominal resolution product. 
% For applications where the NRCS is relevant (e.g. surface stress or wind, oil-spill detection), a low-speckle spatially averaged product including polarimetric information, and a compensation of the known part of the angle of incidence dependence is preferable to a full-resolution product which first needs to be properly despeckled, noting that speckle filtering requires a level of expertise that makes the product deviate from the ARD rationale.

% It also is worth noting that a full-resolution \acs{orb} product is, in general, not a good starting point for the generation of enhanced ocean products. In many cases, the generation of these products requires SLCs (or even raw data), for example for sub-look generation to enable the computation of intensity cross-spectra, or for the optimal suppression of azimuth ambiguities in coastal areas by optimal selection of the azimuth-wavenumber (AKA Doppler frequency) weighting.


% The \acs{orb} product, either as defined in \cite{CEOS-ARD_PFS_ORB} or in the variants discussed in the proposal for this study, can be relatively straightforwardly derived from the \acs{nrb}.

% While many references can be found to prototype implementations of \acs{nrb} products for both Sentinel-1 and ROSE-L, there is limited evidence of similar efforts for \acs{orb}, although \acs{orb}-type products are likely generated as intermediate products by many stakeholders in their analysis of ocean data.

% Lower resolution ocean \textit{roughness} products derived from Sentinel-1 data are, in contrast, provided (for example) through the through Oceandatalab's Ocean Virtual Lab portal.

% \begin{longtblr}[
%   label = tab:orb_products,
%   entry = none,
%   caption = {Data layers of the \acs{orb} },
% ]{
%   width = \linewidth,
%   colspec = {Q[c,20]Q[c,30]Q[l,40]},
%   hlines,
%   hline{1-Z} = {-}{0.05em},
% }
% \textbf{\acs{ard} Product} & \textbf{Description} & \textbf{Notes}\\
% $\sigma_{0,1}$ & Full resolution (single look) radar intensity, in $\sigma_0$ convention & This is the default \acs{orb} product defined in  \cite{CEOS-ARD_PFS_ORB} \\
% $\sigma_{0,L}$ & Spatially-averaged, LxL m$^2$ resolution radar intensity, in $\sigma_0$ convention. & The possibility of multilooked product is implicit in the \acs{ceos} \acs{orb} definition, as it defines metadata regarding speckle filtering, etc. \\
% $\sigma_{0,L}/\sigma_{0,\mathrm{mod}}$ & (Mean) wind normalized backscatter measurements & This product is not required, but it defined in \cite{CEOS-ARD_PFS_ORB} and highly useful as it removes the geometry dependence of the backscatter intensity, allowing a direct interpretation of atmospheric phenomena. The complementary metadata also allow a low-effort wind-retrieval. //
% 1 & $C_{3mod}$(1,1) & HH backscatter & Real  \\  
% 2 & $C_{3mod}$(1,2) & HH·HV* correlation & Complex  \\
% 3 & $C_{3mod}$(1,3) & HH·VV* correlation & Complex \\
% 4 & $C_{3mod}$(2,2) & HV backscatter & Real \\
% 5 & $C_{3mod}$(2,3) & HV·VV* correlation & Complex \\
% 6 & $C_{3mod}$(3,3) & VV backscatter & Real \\
%\end{longtblr}

\subsubsection{Input and Auxiliary Data}
\label{sec:aux_data_ocean}
% To generate the product the following auxiliary data are necessary:
% \begin{itemize}
%     \item external Geoid
%     \item \acs{slc} precise orbits
%     \item \acs{etad}-like product for atmospheric delay and timing corrections.
%     \item \acs{slc} ancillary data
% \end{itemize}

To generate the \acs{orb} product coincide with the \acs{nrb} product, the following input and auxiliary data are necessary: 
\begin{longtblr}[
	label = tab:in_aux_orb_products,
	entry = none,
	caption = {\acs{orb} input and auxiliary data},
	]{
		width = \linewidth,
		colspec = {Q[l,20]Q[l,60]Q[c,20]},
		hlines,
		hline{1-Z} = {-}{0.05em},
	}
	\textbf{Name}   & \textbf{Description}   & \textbf{Mandatory (M) / Optional (O)} \\
	\acs{slc} or \acs{cslc} & L1 Single Look Complex image or co-registered L1 Single Look Complex image and annotations & M \\
    \acs{dem} & External DEM (e.g., Copernicus DEM at different resolutions, with 30 m as the default worldwide). & O \\
    Orbits & \acs{slc} Standard, restituted or precise orbits & M \\
    \acs{etad} & \acs{etad}-like products for atmospheric delay and timing corrections & O \\
    \acs{tec} & External \acs{tec} data for delay correction & O \\
    Geoid & External geoid & M \\
    \acs{era5} & ECMWF \acs{era5} wind data for angle of incidence compensation & M \\
\end{longtblr}

\subsubsection{Main Variables}
\label{sec:variables_orb}
The main variables of the \acs{orb} product are listed as follows:
\begin{longtblr}[
	label = tab:var_orb_products,
	entry = none,
	caption = {\acs{orb} main variables},
	]{
		width = \linewidth,
		colspec = {Q[l,15]Q[l,52]Q[l,20]Q[l,13]},
		hlines,
		hline{1-Z} = {-}{0.05em},
	}
	\textbf{Name} & \textbf{Description} & \textbf{Dimensions} & \textbf{Data Type} \\
    $\sigma_{0,\mathrm{raw}}$ & Full resolution (single look) radar intensity, in \gls{sigma0} convention & polarization, acquisition\_date, x, y & float32 \\
    $\sigma_{0,\mathrm{nesz}}$ & Spatially-averaged, LxL m$^2$ resolution noise equivalent sigma zero, in \gls{sigma0} convention & polarization, acquisition\_date, x, y & float32 \\
    $\sigma_{0,\mathrm{denoise}}$ & Spatially-averaged, LxL m$^2$ resolution denoized radar intensity, in \gls{sigma0} convention & polarization, acquisition\_date, x, y & float32 \\
    mask & Full resolution (single look) bright target mask & polarization, acquisition\_date, x, y & bool \\
    $\sigma_{0,L}$ & Spatially-averaged, LxL m$^2$ resolution denoized and outlier removed radar intensity, in \gls{sigma0} convention & polarization, acquisition\_date, x, y & float32 \\
    $\sigma_{0,L}/\sigma_{0,\mathrm{mod}}$ & Spatially-averaged, LxL m$^2$ resolution (mean) wind normalized backscatter measurements & polarization, acquisition\_date, x, y & float32 \\
\end{longtblr}
The table below describes the dimensions/coordinates associated with the \acs{orb} main variables.
\begin{longtblr}[
	label = tab:var_orb_coordinates,
	entry = none,
	caption = {\acs{orb} main variables coordinates/dimensions},
	]{
		width = \linewidth,
		colspec = {Q[l,20]Q[l,62]Q[l,18]},
		hlines,
		hline{1-Z} = {-}{0.05em},
	}
	\textbf{Name} & \textbf{Description} & \textbf{Data Type} \\
	polarization & Acquired polarizaions & string \\
    acquisition\_date & Image acquisition dates, formatted as YYYY-MM-DD & string \\
    x & UTM/UPS x (East) coordinates [m] & float64 \\
    y & UTM/UPS y (North) coordinates [m] & float64 \\
\end{longtblr}
\subsubsection{Main Metadata}
\label{sec:orb_product_metadata}
The following table lists the main metadata and per-pixel metadata associated with the \acs{orb} product.

\begin{longtblr}[
	label = tab:meta_orb_products,
	entry = none,
	caption = {\acs{orb} main metadata and per-pixel metadata},
	]{
		width = \linewidth,
		colspec = {Q[l,17]Q[l,38]Q[l,20]Q[l,13]Q[c,12]},
		hlines,
		hline{1-Z} = {-}{0.05em},
	}
	\textbf{Name} & \textbf{Description} & \textbf{Dimensions} & \textbf{Data Type} & \textbf{Per-pixel Metadata}\\
	reference\_slc / reference\_cslc & List of
    links to input reference \acs{slc} or \acs{cslc} & - & list of string & \\
    reference\_geoid & Link to reference GEOID. Only for \acs{eorb} & - & string & \\
    tec & Link to external \acs{tec} data, if used. & - & string & \\ 
    era5 & Link to external \acs{era5} data, if used.  & - & string & \\
        x\_sample & Pixel spacing in x direction & - & float32 &\\
        y\_sample & Pixel spacing in y direction & - & float32 &\\
        ground\_heading & Azimuth direction & - & float32 & \\
        look\_direction & Planar angle between	north and each range direction & - & float32 &\\
        inc\_angle & Angle of incidence & x, y & float32 & \ding{51} \\
        nlooks & Number of equivalent looks & x, y & float32 & \ding{51} \\
        azimuth\_time & Azimuth time & x, y & datetime64 & \ding{51} \\
        slant\_range\_time & Slant range time & x, y & datetime64 & \ding{51} \\
\end{longtblr}

The table below describes the dimensions/coordinates associated with the \acs{orb} main metadata.

\begin{longtblr}[
	label = tab:var_orb_metadata_coordinates,
	entry = none,
	caption = {\acs{orb} main metadata coordinates/dimensions},
	]{
		width = \linewidth,
		colspec = {Q[l,20]Q[l,62]Q[l,18]},
		hlines,
		hline{1-Z} = {-}{0.05em},
	}
	\textbf{Name} & \textbf{Description} & \textbf{Data Type} \\
	x & UTM/UPS x (East) coordinates [m] & float64 \\
    y & UTM/UPS y (North) coordinates [m] & float64 \\
\end{longtblr}

% The indicated main metadata for these products derive from the analysis of \acs{ceos} Product Format Specification \cite{CEOS-ARD_PFS_SAR}.% and of the metadata available information for the products list in \autoref{tab:orb_products}.

% Most of the listed products are declared compliant with CEOS requirements whereby it is assumed that their metadata include at least the threshold requirements of \cite{CEOS-ARD_PFS_ORB} and \cite{CEOS-ARD_PFS_SAR}.

% In \autoref{tab:metadata_orb_products}, ``All" means that the correspondent metadata are present in all the products in \autoref{tab:orb_products} and in CEOS recommendations. %Note that NISAR \acs{rtc} is included in the NISAR \acs{gcov} product. 
% This table can be understood as a \acs{orb}-specific revision of \autoref{tab:metadata_nrb_products}.

% \begin{longtblr}[
%   label = tab:metadata_orb_products,
%   entry = none,
%   caption = {ORB product main metadata from \acs{sota} analysis},
% ]{
%   width = \linewidth,
%   colspec = {Q[l,15]Q[l,42]Q[l,18]Q[l,25]},
%   hlines,
%   hline{1-Z} = {-}{0.05em},
%   rowhead=1,
% }
% \textbf{Metadata name}   & \textbf{Description}   & \textbf{Reference Product} & \textbf{Notes} \\
% Product Type & Product type name & All & - \\
% Source data attributes & Subsection describing \acs{sar} acquisitions used to generate the product. Main attributes for each source are:

% -Source Data Access

% -Instrument

% - Acquisition Time

% - Acquisition Parameters (Band, Centre frequency, Observation mode, Polarization(s), Antenna pointing, Beam ID)

% - Orbit Information

% - Processing Parameters

% - Image Attributes (geometry, azimuth and range pixel spacing and resolution, incidence angles, range distance from the sensor to the surface)

% & All & - \\
% Product Data Access & Processing parameters details of the product (processing facility, processing date, software version, \acs{url} or \acs{doi} from where product can be retrieved.
%  & All & - \\
% Polarization(s) & Indication of the polarization. & All & - \\
% Product Sample Spacing & - Pixel (column) spacing

% -	Line (row) spacing
%  & All & - \\
% Product Equivalent Number of Looks & Equivalent Number of Looks (\acs{enl}) &  & Not required by \acs{ceos} (listed as ”goal”). 
%  \\
% Product Resolution & Average spatial resolution of the product along columns and rows
%  & - & Not required by \acs{ceos} (listed as ”goal”).\\

% Range of validity & Variable range of validity &  & The analysed documentation does not explicitly mention the presence of this field for all the other indicated products. \\
% Product Filtering & Flag indicating if speckle or noise filter has been applied. Metadata should include reference to algorithm as \acs{doi} or \acs{url} and input filtering parameters. & All & As threshold. Not relevant in \acs{ceos} \acs{orb} specification, but necessary if we consider, as strongly recommended, lower resolution speckle-filtered products. \\
% Product Bounding Box & Two opposite corners of the product file (bounding box, including any zero-fill values) are identified, expressed in the coordinate reference system.
% Four corners of the product file are recommended for scenes crossing the Antimeridian, or the North or the South Pole. 
%  & All & - \\
% Product Geographical Extent & The geometry of the \acs{sar} image footprint expressed in WGS84, in a standardised format (e.g., WKT Polygon). & All & - \\
% Product Image Size & 
% - Number of lines

% - Number of pixels per line

% - File header size (if applicable)

% - Number of no-data border pixels (if appl.)
%  & All & - \\
% Backscatter Measurements & "Geoid-corrected"  Sigma-Nought ($\sigma^0$) convention, for each polarization.
% Metadata should include:

% - Measurement Type: Sigma-Nought (NRCS)

% - Backscatter Expression Convention: Linear amplitude or linear power

% - Polarization: HH, HV, VV, VH

% - Data Format: GeoTIFF, HDF5, NetCDF, etc.

% - Data Type: Integer, Float, etc.

% - Bits per Sample

% - Byte Order

% & CEOS \acs{orb} & Partially present in all the listed products. \\
% Scaling Conversion & Equation to convert linear amplitude/power to logarithmic decibel scale, including, if applicable, the associated calibration (dB offset) factor, and/or the equation used to convert compressed data (int8/int16/float16) to float32. &  CEOS \acs{orb} & - \\

% GMF & Geophysical model function used to relate model winds to reference NRCS, $\sigma_{0,\mathrm{mod}}$. & & This could be provided as a mathematical model with or as a reference and URL to a GMF description and implementation. \\

% FR method & Method used to estimate FR. & &  Under some assumptions, FR can be estimated from data, or, alternatively, from ionospheric models. \\
% %
% % \acs{rtc} Algorithm & Reference to \acs{rtc} algorithm and sources of used auxiliary data. & All & - \\
% % Product Pixel Coordinate Convention & Coordinate referring to the Centre, or the Upper Left Corner or the Lower Left Corner of a pixel & All & - \\
% Product Coordinate Reference System & Parameters required to geolocate data in that map projection, expressed in a standardised format (e.g., WKT). 

% Indicate \acs{epsg} code, if defined for the \acs{crs}
%  & All & - \\
% Reference Orbit & Absolute orbit number of reference \acs{slc}. It includes the baseline for the source data relative to the reference orbit used (for approximate use only).
%  & All & The generation of \acs{orb}s is on single \acs{slc}. \\
% Geometric Correction Algorithms & Information on the algorithm used for geometric correction & - & Not required by \acs{ceos} \\
% % Digital Elevation Model & - Reference to Digital Elevation Model used for co-registration step, ortho-rectification and \acs{rtc} procedure.

% % - Reference to Earth Gravitational Model (\acs{egm}) used for geometric correction.
% %  & All & All steps, requiring it, use the same \acs{dem}.\\
% Radiometric Accuracy & Uncertainty on  $\sigma^0$. & - & Not required by \acs{ceos} (listed as ”goal”). \\
% Geometric Accuracy & Geolocation accuracy estimates, based on SLC assessment (\acs{slc} \acs{ale}) and/or measured directly in the final projected final product (\acs{ard} \acs{ale}). &  & Not required by \acs{ceos} (listed as ”goal”). 

% %For QA4EO RTC, the geometric accuracy does not include the relative RMSE and reference.
% \\
% Gridding Convention & All necessary information of used grid and used convention & All & - \\
% Processing information & Information on the processing: common band filtering, spatial averaging method, \acs{rfi} correction, co-registration algorithm, etc.. & All & Some of this information could be already available through source-data
% pointers. \\
% Units & Units of the variables of measurements. & All & Not found in \acs{ceos}.
% %\centering
% \end{longtblr}

% \autoref{tab:per-pixel_metadata_orb_products} indicates the main product per-pixel metadata. The names and definitions of the majority of the per-pixel metadata are exacted from \acs{ceos} Product Format Specification \cite{CEOS-ARD_PFS_SAR}.% and of the per-pixel metadata available information for the products list in \autoref{tab:orb_products}.

% Most of the listed products are declared compliant with CEOS requirements whereby it is assumed that their per-pixel metadata includes at least the threshold requirements of \cite{CEOS-ARD_PFS_ORB} and \cite{CEOS-ARD_PFS_SAR}.

% In \autoref{tab:per-pixel_metadata_orb_products}, "All" means that the correspondent per-pixel metadata is present in all the products in \autoref{tab:orb_products} and in CEOS recommendations. %Note that NISAR \acs{rtc} is included in the NISAR \acs{gcov} product.

% \begin{longtblr}[
%   label = tab:per-pixel_metadata_orb_products,
%   entry = none,
%   caption = {\acs{orb} product main per-pixel metadata from \acs{sota} analysis},
% ]{
%   width = \linewidth,
%   colspec = {Q[l,15]Q[l,33]Q[l,18]Q[l,34]},
%   hlines,
%   hline{1-Z} = {-}{0.05em},
% }
% \textbf{Metadata name}   & \textbf{Description}   & \textbf{Reference Product} & \textbf{Notes} \\
% Data mask & Mask image indicating:

% - Valid data

% - Invalid data

% - No data & All & - \\
% Equivalent Number of Looks Image & Equivalent Number of Looks (\acs{enl}) &  & Not required by \acs{ceos}, but important for characterization of radiometric quality for optionally spatially-averaged products. \\
% Scattering Area & geoid-based scattering area image &  & Not required by \acs{ceos} (indicated as a goal requirement). It depends on the used algorithm, e.g. if number of independent input pixels is constant this areas is going to vary for each output pixel. \\
% Local incidence angle & DEM based local incidence angle image & All &  \\
% Ellipsoidal incidence angle & Ellipsoidal incidence angle image & %NISAR RTC (GCOV), QA4EO RTC & Not required by \acs{ceos} (indicated as goal requirement). The analysed ROSE-L documentation does not explicitly mention the presence of this field.
% \\
% % Per-pixel \acs{dem} and per-pixel GEOID & DEM and/or GEOID using for the processing & CEOS \acs{nrb} & This can be useful for a variety of operations (undoing the geometric correction, the atmospheric correction, etc.). Probably it is not necessary report the DEM/GEOID for each pixel of the product: it could be sufficient the information from "Digital Elevation Model" metadata. \\
% Noise Power Image & Estimated noise equivalent $\sigma^0$  & CEOS \acs{orb} & - \\
% % Gamma-to-Sigma Ratio Image & Ratio between the integrated area in the Gamma projection and the integrated area in the Sigma projection (on the ground) & All & - \\
% Radiometric Accuracy & Estimated uncertainty (e.g. $1\sigma$) of the intensity products. & CEOS \acs{orb} & Not required by CEOS but important for some applications, for example for uncertainty propagation to surface winds estimates \\
% $U_{10s,m}$ & Stress-equivalent model 10-m wind used to generate reference intensity, $\sigma_{0,\mathrm{mod}}$  & - & -\\
% Faraday Rotation & Estimated FR per pixel & - & Not specified in CEOS \acs{orb}, but potentially relevant to correct or address leakage of co-polar signal in cross-polar channel.


% %\centering
% \end{longtblr}

\subsubsection{Reference System Tiling and Slicing}
A gridding/sampling frame, consistent with the reference system and slicing defined for the \acs{nrb}, is generally used.
\begin{longtblr}[
	label=tab:orb_ref_system_slicing,
	entry = none,
	caption = {Reference system and Slicing for \acs{orb}},
	]{
		width=\linewidth,
		colspec = {Q[l,15]Q[l,35]Q[l,35]Q[l,15]},
		%vline{2-5} = {-}{},
		%vline{5-5} = {-}{},
		hlines,
		hline{1-Z} = {-}{0.05em},
        rowhead=2
	}
	\textbf{Product} & \textbf{Reference System} & \textbf{Slicing} & \textbf{Spacing}\\
	\textbf{\acs{orb}} (default) & Cartographic: UTM/UPS & US-\acs{mgrs} 100 km x 100 km & 10 m\\
    \textbf{\acs{orb}} (optional) & Cartographic: UTM/UPS & US-\acs{mgrs} 100 km x 100 km & 30 m\\	
\end{longtblr}
\subsection{Product Requirements}
\subsubsection{Accuracy}
Main impacts on \acs{orb} product generation and main requirements are similar to those indicated in the correspondent section of the \acs{nrb} product. The limitations, in this case, are less stringent for the local slope (here the geoid is used) than for the geolocation accuracy (less binding over ocean).
\begin{longtblr}[
	label = tab:orb_accuracy_cband,
	entry = none,
	caption = {Accuracy Requirements for \acs{orb} product (C-Band)},
	]{
		width = \linewidth,
		colspec = {Q[l,30] Q[l,60] Q[c,20]},
		hlines,
		hline{1-Z} = {-}{0.05em},
		rowhead = 1,
	}
	
	\textbf{Metric} & \textbf{Description} & \textbf{Requirement} \\
	
	\SetCell[r=2]{l} Range Accuracy (\%) 
	& Model-Driven (only orbit and \acs{dem}) & $\leq 5\%$\\
	& Model-Driven + atmospheric correction & $\leq 2\%$ \\
	
	\SetCell[r=2]{l} Azimuth Accuracy (\%) 
	& Model-Driven (only orbit and \acs{dem}) & $\leq 2\%$\\
	& Data-Driven (\acs{esd}) & $\leq 0.5\%$ \\
	
	\acs{rcs} difference (dB) & Stability of the \acs{rcs} comparing different acquisition geometries over the same area considering a resolution of 30m & $\leq 5$~dB \\
	
\end{longtblr}

\begin{longtblr}[
	label = tab:orb_accuracy_lband,
	entry = none,
	caption = {Accuracy Requirements for \acs{orb} product (L-Band)},
	]{
		width = \linewidth,
		colspec = {Q[l,30] Q[l,60] Q[c,20]},
		hlines,
		hline{1-Z} = {-}{0.05em},
		rowhead = 1,
	}
	
	\textbf{Metric} & \textbf{Description} & \textbf{Requirement} \\
	
	\SetCell[r=3]{l} Range Accuracy (\%) 
	& Model-Driven (only orbit and \acs{dem}) & $\leq 60\%$\\
	& Model-Driven + atmospheric correction & $\leq 25\%$ \\
	& Model-Driven + atmospheric correction + Data-Driven (\acs{icc}) & $\leq 10\%$ \\
	
	\SetCell[r=3]{l} Azimuth Accuracy (\%) 
	& Model-Driven (only orbit and \acs{dem}) & $\leq 150\%$ \\
    & Model-Driven + atmospheric correction & $\leq 50\%$ \\
	& Data-Driven (\acs{icc}) & $\leq 10\%$ \\
	
	\acs{rcs} difference (dB) & Stability of the \acs{rcs} comparing different acquisition geometries over the same area considering a resolution of 30m & $\leq 5$~dB \\
	
\end{longtblr}

\subsubsection{Timeliness}
For many potential use cases of \acs{orb} products, timeliness can be critical, as the information in the image is highly ephemeral. The production of the product has the same timeline of \acs{nrb} product. In case of \acs{orb} being the requirement on geolocation more relaxed, it is possible to work only with nominal geometry without need for refinement of timing. Atmospheric corrections to produce unbiased NRCS estimates and to provide ancillary surface wind estimates are needed to allow the use of the \acs{orb} product for atmospheric applications (wind estimation, Marine Atmospheric Boundary Layer (MABL) parameterization, etc). \autoref{fig:orb_single} shows the flowchart of generating the \acs{orb} product using a single \acs{slc} image.
\begin{figure}[H]
	\centering
	\resizebox{\textwidth}{!}{
		\begin{tikzpicture}[node distance=0.5cm and 2cm]
			
			% Define the style for the nodes
			% \tikzstyle{node} = [circle, draw, fill=blue!20, text centered, minimum size=1.5cm]
			
			% Draw nodes
			\node[rectin] (init) {\acs{slc}};
			\node[block, right=of init] (g2r) {Geo-to-Radar};
			\node[block, right=of g2r] (rtc) {\acs{rtc}};
			\node[block, right=of rtc] (geocoding) {Geocoding};
                \node[block, right=of geocoding] (denoise) {Denoise};
                \node[block, below=of denoise] (outlier) {Outlier removal};
                \node[block, below=of outlier] (downsample) {Downsample};
                \node[block, below=of downsample] (roughness) {Normalization};
                \node[rectout, below=of roughness] (orb) {\acs{orb}};
			\node[rectaux, below=of init] (dem) {\acs{dem}};
                \node[empty, below=of g2r] (empty0) {};
			\node[rectaux, below=of empty0] (orbit) {Orbital data};
                \node[rectaux, left=of downsample] (era5) {ERA5 wind};
			\node[rectauxd, below=of orbit] (models) {External models};
                \node[rectauxd, right=of downsample] (cmod5n) {CMOD5N};
			% \node[rectind, below=of models] (ref_stack) {Reference SLC stack};
			
			% Draw edges with arrows
			\draw[->, thick] (init) -- (g2r);
			\draw[->, thick] (g2r) -- (rtc);
			\draw[->, thick] (rtc) -- (geocoding);
			\draw[->, thick] (dem) -| (g2r);
			\draw[->, thick] (orbit) -- (g2r);
			\draw[dotted, thick] (models) -- (orbit);
			% \draw[->, thick] (ref_stack) -| (g2r);
			\draw[->, thick] (geocoding) -- (denoise);
                \draw[->, thick] (denoise) -- (outlier);
                \draw[->, thick] (outlier) -- (downsample);
                \draw[->, thick] (downsample) -- (roughness);
                \draw[->, thick] (roughness) -- (orb);
                \draw[->, thick] (era5) -- (downsample);
                \draw[->, thick] (cmod5n) -- (downsample);
                % Timeline with linked markers and adaptive length
			\def\timelineStart{init}
			\def\timelineEnd{orb}
			\draw[thick, ->] (\timelineStart.west |- 0,-8) -- (\timelineEnd.east |- 0,-8) node[right] {Time};
			\foreach \node/\label in {dem/Start, models/\textbf{1-30d}, orb/<10min} {
				\draw (\node.south |- 0,-8) -- (\node.south |- 0,-8.5) node[below] {\label};
				\draw[dotted, thick] (\node.south) -- (\node.south |- 0, -8.25); % Connects the node to the timeline
			}
			
		\end{tikzpicture}
	}
	\caption{Sample workflow for generation of \acs{orb} considering each \acs{slc} image separately.}
	\label{fig:orb_single}
\end{figure}


\subsection{Usability, Integrability and Interoperability}
\label{sec:orb_usability_integrability}
The characteristics of the \acs{orb} products in terms of usability, integrability and interoperability  are in the following table.

\begin{longtblr}[
  caption = {Usability, Integrability, and Interoperability of \acs{orb} product},
  label = {tab:orb_usability_integrability_interoperability},
]{
  colspec = {Q[2cm] Q[12cm]},
  rowhead = 1,
  width = \linewidth
}
\hline
\textbf{Aspect} & \textbf{Description} \\
\hline
Usability &
\begin{itemize}
    \item Pre-Calibrated Data: the \acs{orb} product is radiometrically calibrated, meaning users do not need to perform additional steps to correct the radar backscatter values. This simplifies the analysis for the users, including those who may not be experts in radar data processing.
    \item Normalized Backscatter: since the backscatter values are normalized (converted into \gls{sigma0} by applying a geoid correction procedure), they can be easily compared across different acquisitions, sensors, or times. 
    \item Geometrically Corrected (Georeferenced): \acs{orb} products are orthorectified and projected into standard map coordinates, reducing the need for users to handle raw slant-range data or worry about terrain-induced distortions (unless further specific geometric corrections are needed for the user's application).
    \item Scalable and Flexible Data Formats: \acs{orb} products are typically provided in common geospatial formats (such as GeoTIFF, NetCDF, etc.), which are compatible with a wide range of geospatial tools (e.g., QGIS) and programming libraries (e.g., GDAL, Python). This facilitates easy integration into existing analysis workflows.
    \item Ready for Machine Learning: \acs{orb} products, being standardized and normalized, can be directly ingested into machine learning models for classification, regression, or segmentation tasks without the need for extensive preprocessing. Their calibrated and consistent nature makes them suitable for training AI models on large datasets.
\end{itemize}
\\
Integrability &
\begin{itemize}
    \item Integration into \acs{gis} platforms: \acs{orb} products are easy to integrate into Geographic Information Systems (\acs{gis}) due to their georeferencing and common data formats. Users can overlay these products with vector layers, other remote sensing data (optical, infrared), or ground truth data for comprehensive spatial analysis.
    \item Cloud-based processing: \acs{orb} products can be easily integrated into cloud-based processing environments like Google Earth Engine (\acs{gee}), \acs{aws}, or Sentinel Hub, etc. This allows large-scale, on-demand processing without requiring local resources. Users can perform global or regional analysis by leveraging cloud computing for processing extensive datasets.
    \item Compatibility with remote sensing software: software packages like ESA SNAP (Sentinel Application Platform), ENVI, or commercial and free GIS software like ArcGIS and QGIS can directly process and analyse \acs{orb} products, making them easy to integrate into radar data analysis pipelines.
\end{itemize}
\\
Interoperability &
\begin{itemize}
  \item Cross-mission compatibility for ROSE-L and Sentinel-1 missions
  \item Variety of applications, particularly for monitoring and analysing oceanic and coastal environments: Sea Surface Wind Monitoring, Wave Height and Ocean State Estimation, Ocean Current Mapping, Oil Spill Detection, Sea Ice and Iceberg Monitoring, Coastal Erosion and Sediment Transport Analysis, Ship Detection and Maritime Surveillance.
  \item High interoperability due to standard geographic referencing and pixel alignment across scenes
\end{itemize}
\\
\hline
\end{longtblr}

% \begin{itemize}
%     \item Usability.
%         \begin{itemize}
%             \item Pre-Calibrated Data: the \acs{orb} product is radiometrically calibrated, meaning users do not need to perform additional steps to correct the radar backscatter values. This simplifies the analysis for the users, including those who may not be experts in radar data processing.
%             \item Normalized Backscatter: since the backscatter values are normalized (converted into sigma-nought by applying a geoid correction procedure), they can be easily compared across different acquisitions, sensors, or times. %This makes it highly useful for time series analysis, change detection, and long-term monitoring without worrying about differences in acquisition geometry or sensor calibration. This increases the product’s usability for land cover classification, land use analysis, and surface change monitoring.
%             \item Geometrically Corrected (Georeferenced): \acs{orb} products are orthorectified and projected into standard map coordinates, reducing the need for users to handle raw slant-range data or worry about terrain-induced distortions (unless further specific geometric corrections are needed for the user's application).
%             \item Scalable and Flexible Data Formats: \acs{orb} products are typically provided in common geospatial formats (such as GeoTIFF, NetCDF, etc.), which are compatible with a wide range of geospatial tools (e.g., QGIS) and programming libraries (e.g., GDAL, Python). This facilitates easy integration into existing analysis workflows.
%             \item Ready for Machine Learning: \acs{orb} products, being standardized and normalized, can be directly ingested into machine learning models for classification, regression, or segmentation tasks without the need for extensive preprocessing. Their calibrated and consistent nature makes them suitable for training AI models on large datasets.
%         \end{itemize}
%     \item Integrability into different systems and workflows.
%         \begin{itemize}
%             \item Integration into \acs{gis} platforms: \acs{orb} products are easy to integrate into Geographic Information Systems (\acs{gis}) due to their georeferencing and common data formats. Users can overlay these products with vector layers, other remote sensing data (optical, infrared), or ground truth data for comprehensive spatial analysis.
%             % \item Interoperability with other \acs{ard} products: the \acs{nrb} could be integrated with other \acs{ard} products (such as optical data, LiDAR, or \acs{sar} products) for multi-sensor analysis. For instance, combining \acs{nrb} data with optical satellite imagery can enhance land cover classification or environmental monitoring efforts. In addition, they can complement \acs{sar} interferometric products (e.g., interferograms, coherence maps) by providing context on surface properties and changes in scattering characteristics. This is valuable in areas like infrastructure monitoring, natural hazard detection, and land subsidence studies.
%             \item Cloud-based processing: \acs{orb} products can be easily integrated into cloud-based processing environments like Google Earth Engine (\acs{gee}), \acs{aws}, or Sentinel Hub, etc.. This allows large-scale, on-demand processing without requiring local resources. Users can perform global or regional analysis by leveraging cloud computing for processing extensive datasets.
%             \item Compatibility with remote sensing software: software packages like ESA SNAP (Sentinel Application Platform), ENVI, or commercial and fre GIS software like ArcGIS and QGIS can directly process and analyse \acs{orb} products, making them easy to integrate into radar data analysis pipelines.
%             % \item Environmental and Agricultural Monitoring Systems: \acs{nrb} products can be integrated into environmental monitoring platforms to assess parameters like soil moisture, forest cover, and agricultural productivity. By coupling radar data with models or indices, they offer insights into ecosystem health or crop conditions even under cloud cover or adverse weather conditions.
%         \end{itemize}
% \end{itemize}

% % The usability of NRB products is highly dependent on the spatial and temporal resolution. For high-frequency monitoring or fine-scale analysis, users must ensure that the \acs{nrb} product offers the necessary resolution. For example, low-resolution \acs{nrb} products might be unsuitable for fine-scale monitoring.

% % \acs{nrb} products often include multiple polarizations (e.g., VV, VH), which provide complementary information about surface properties. The ability to use different polarizations enhances integrability into diverse applications like vegetation analysis (VH polarization is sensitive to biomass) or water body detection (VV polarization).

% % In case of integration with other products, it could be necessary and easier to define a common data format, reference system, and geographic grid.

% % In some of the analysed \acs{nrb} products the speckle filter is not applied. For specific analyses like land cover classification or vegetation monitoring, speckle noise mitigation may be required. This is an additional preprocessing step that one may need to consider.

% % Finally, radar backscatter can be influenced by surface roughness, moisture content, and other environmental factors. These must be accounted for when interpreting the data, especially in applications where surface characteristics (e.g., soil moisture, vegetation growth) vary over time.

% The Ocean Radar Backscatter (\acs{orb}) product has a variety of applications, particularly for monitoring and analysing oceanic and coastal environments. Key applications include:
% \begin{itemize}
%     \item Sea Surface Wind Monitoring,
%     \item Wave Height and Ocean State Estimation,
%     \item Ocean Current Mapping,
%     \item Oil Spill Detection,
%     \item Sea Ice and Iceberg Monitoring,
%     \item Coastal Erosion and Sediment Transport Analysis,
%     \item Ship Detection and Maritime Surveillance.
% \end{itemize}

% % Normalized Radar Backscatter products are versatile for a wide range of applications, from environmental monitoring to infrastructure management, disaster response, and resource exploration. Below are some of the key applications of \acs{orb} products:
% % \begin{itemize}
% %     \item Land Cover Classification.
% %         \begin{itemize}
% %             \item Vegetation Monitoring: \acs{orb} products can help in identifying different types of vegetation and monitoring their health, structure, and density, as different land covers (e.g., forests, agricultural fields, grasslands) have distinct radar backscattering signatures.
% %             \item Urban vs. Rural Classification: Radar backscatter can distinguish between built-up areas and natural land covers.
% %         \end{itemize}
% %     \item Agriculture.
% %         \begin{itemize}
% %             \item Crop Monitoring and Yield Prediction: \acs{orb} can be used to monitor crop growth stages and predict yields by identifying differences in crop structures.
% %             \item Soil Moisture Estimation: SAR backscatter is sensitive to soil moisture content, making \acs{orb} products useful for assessing soil moisture in agricultural fields, which is crucial for water management and irrigation practices.
% %         \end{itemize}
% %     \item Flood Detection: Radar backscatter is affected by water, enabling \acs{orb} products to detect flooded areas
% %         \begin{itemize}
% %             \item in agricultural lands after heavy rains or during irrigation events.
% %             \item in urban zones
% %             \item for tracking water levels in real-time, even during cloudy or rainy conditions
% %             \item for creating flood risk maps
% %         \end{itemize}
% %     \item Change Detection
% %         \begin{itemize}
% %             \item for disaster management (earthquake and post-disaster damage assessment, forest fire and burned area mapping)
% %             \item for deforestation monitoring
% %             \item for urbanization monitoring
% %             \item for volcanic monitoring (e.g., appearance of new lava flows)
% %             \item for monitoring the polar ice caps
% %         \end{itemize}
% %     \item Forest Biomass Estimation \& Carbon Stock Estimation for climate change monitoring
% %     \item Water Resource Management.
% %         \begin{itemize}
% %             \item Wetland Mapping: The ability of \acs{orb} to differentiate between water and land makes it useful for mapping and monitoring wetlands and other water bodies.
% %             \item River and Lake Monitoring: \acs{orb} products can help monitor changes in river flows and lake levels, essential for hydrology studies and water resource management.
% %         \end{itemize}
% %     \item Military and Security Applications: \acs{orb} products can detect changes in terrain, infrastructure, border and coast or activity in a region, even under poor visibility, making them valuable for military and security operations.
% %     \item Environmental Monitoring.
% %         \begin{itemize}
% %             \item Coastal Erosion Monitoring: \acs{orb} products are useful in tracking shoreline changes and coastal erosion, which is important for coastal management and protection.
% %             \item Desertification and Land Degradation: to monitor areas undergoing desertification or other forms of land degradation by detecting changes in surface properties.
% %             \item Monitoring of Environmental Protection Areas: \acs{orb} products can be used to track illegal activities, such as deforestation or construction, in protected areas.
% %         \end{itemize}
% % \end{itemize}
